\clearpage
\thispagestyle{empty}
\pdfbookmark[0]{Ausgabe und Quellen}{translation-sources}
\begin{center}\textbf{Ausgabe und Quellen}\end{center}
\noindent Diese deutsche Fassung enthält die historischen Titelblätter, das Vorwort, das Inhaltsverzeichnis und alle 192 Artikel aus Carl Friedrich Gauß’ \textit{Theoria motus corporum coelestium in sectionibus conicis solem ambientium} (Hamburg, 1809) nach dem Abdruck in seinen \textit{Werken}, Band~VII (Göttingen und Leipzig, 1906), einschließlich der drei Tafeln, der vier Abbildungen und Brendels historischer Bemerkungen. Sie enthält außerdem alle 18 Zusätze zur Theoria motus und die zugehörigen historischen Bemerkungen. Es folgen die neun Notizen und Briefe zur Bestimmung von Kometenbahnen sowie Gauß’ neue Tafel für die parabolische Bewegung mit vollständiger Erläuterung und allen drei Teilen. Die Briefe an Encke, Olbers und den Herausgeber sowie Brendels Bemerkungen schließen den Teil zur parabolischen Bewegung ab. Der Teil über die Störungen der Ceres enthält beide Methoden, die Briefe und Tafeln sowie Brendels Bemerkungen. Der Teil über die Störungen der Pallas enthält alle 37 Briefe der Ausgabe von 1906 und die vollständige \textit{Darstellung einer neuen Methode zur Berechnung der Planetenstörungen} mit beiden Entwürfen des vierten Abschnitts. Es folgen beide Rechnungen der speziellen Störungen durch Jupiter, die erste Rechnung der allgemeinen Störungen und die zweite Rechnung bis zum Ende des ersten mit [11] bezeichneten Abschnitts auf S. 561 der Ausgabe von 1906. Der zweite mit [11] bezeichnete Abschnitt auf S. 561–564 liegt außerhalb des Umfangs dieser Ausgabe. Darauf folgen die Tafeln für die Jupiterstörungen sowie die Abschnitte über die Störungen durch Saturn und Mars. Zum Teil über Pallas gehören Brendels Bemerkungen. Die Theorie der Mondbewegung umfasst alle fünf Abschnitte und Brendels Bemerkungen. Den Abschluss bilden die historischen Berichtigungen und Zusätze, F. Kleins Bemerkungen, das historische Inhaltsverzeichnis mit dem Druckvermerk, das unnummerierte Leerblatt und der Verlagskatalog. Die Preisangaben von 25 Mk. (A. IV), 30 Mk. (A. VII) und 180 Mk. (B, Bände I–VI) folgen den Lesungen der schwer erkennbaren Ziffern im Exemplar von 1906. 
\par\medskip\noindent Der historische Text folgt dem vom \href{https://gdz.sub.uni-goettingen.de/id/PPN236008730}{Göttinger Digitalisierungszentrum} digitalisierten Exemplar. Die unsignierte Titelvignette zeigt dasselbe Motiv wie auf dem Titelblatt von 1906. Die Reproduktion stammt aus dem \href{https://commons.wikimedia.org/w/index.php?oldid=1179960952}{zweiten Abdruck von Band I aus dem Jahr 1870}, Exemplar der University of California Libraries, mit der Kennzeichnung \href{https://creativecommons.org/publicdomain/mark/1.0/}{Public Domain Mark 1.0}.
\par\medskip\noindent KI-Übersetzung aus dem Lateinischen und Französischen, KI-Übertragung ins moderne Deutsch und KI-Satz.
